% Minkyu Jeon -- CV (November 2026)
% Positioning: agentic AI + generative models for de novo protein design,
% grounded in cryo-EM and protein conformational heterogeneity.
% Compiles with pdflatex (TeX Live / Overleaf); only standard packages.

\documentclass[10pt,letterpaper]{article}

\usepackage[left=0.7in,right=0.7in,top=0.6in,bottom=0.6in]{geometry}
\usepackage[T1]{fontenc}
\usepackage[utf8]{inputenc}
\usepackage{charter}
\usepackage[scaled=0.92]{helvet}
\usepackage[protrusion=true,expansion=true]{microtype}
\usepackage{xcolor}
\usepackage{enumitem}
\usepackage{tabularx}
\usepackage{array}
\usepackage{hyperref}

% ---------------------------------------------------------------- style
\definecolor{accent}{HTML}{1A4F8B}
\definecolor{muted}{HTML}{5A5A5A}
\definecolor{highlight}{HTML}{B03A2E}
\hypersetup{colorlinks=true,urlcolor=accent,linkcolor=accent,
  pdftitle={Minkyu Jeon - CV},pdfauthor={Minkyu Jeon}}
\urlstyle{same}

\pagestyle{empty}
\setlength{\parindent}{0pt}
\setlength{\parskip}{0pt}
\setlength{\tabcolsep}{3pt}
\setlist[itemize]{leftmargin=1.05em, topsep=1pt, itemsep=1pt, parsep=0pt,
  label={\color{accent}\small\textbullet}}

\newcommand{\me}{\mbox{\textbf{Minkyu Jeon}}}
\newcommand{\eq}{\textsuperscript{*}}
\newcommand{\award}[1]{\textcolor{highlight}{\sffamily\bfseries\small #1}}
\newcommand{\links}[1]{\enspace{\small #1}}
\newcommand{\sep}{\enspace{\color{muted}|}\enspace}

\newcommand{\cvsection}[1]{%
  \vspace{6pt}%
  {\sffamily\bfseries\color{accent}\textls[60]{\MakeUppercase{#1}}}\par
  \vspace{2pt}{\color{accent}\hrule height 0.5pt}\vspace{4pt}}

% \entry{title}{date}{role / subtitle}{location}
\newcommand{\entry}[4]{%
  \begin{tabularx}{\linewidth}{@{}X r@{}}
    {\bfseries #1} & {\sffamily\small\color{muted}#2}\\
    {\itshape #3}  & {\small\color{muted}#4}
  \end{tabularx}\par\vspace{1pt}}

% \entryone{title}{right}: single-line entry (no subtitle row)
\newcommand{\entryone}[2]{%
  \begin{tabularx}{\linewidth}{@{}X r@{}}
    {\bfseries #1} & {\sffamily\small\color{muted}#2}
  \end{tabularx}\par\vspace{1pt}}

% \role{role}{date}: an extra role line under the same organization
\newcommand{\role}[2]{%
  \begin{tabularx}{\linewidth}{@{}X r@{}}
    {\itshape #1} & {\sffamily\small\color{muted}#2}
  \end{tabularx}\par\vspace{1pt}}

% \paper{venue}{award tag}{title}{authors and links}
\newcommand{\paper}[4]{%
  \begin{tabularx}{\linewidth}{@{}>{\raggedright\arraybackslash}p{2.8cm}>{\raggedright\arraybackslash}X@{}}
    \textcolor{accent}{\sffamily\bfseries\small #1} & {\bfseries #3}\\
    \award{#2} & #4
  \end{tabularx}\par\vspace{3pt}}

% two-column list (talks, honors, service, teaching, skills)
% (tabularx inside \newenvironment must use the \tabularx/\endtabularx form)
\newenvironment{rows}[1][2.8cm]%
  {\tabularx{\linewidth}{@{}>{\raggedright\arraybackslash\leavevmode\sffamily\small\color{muted}}p{#1}>{\raggedright\arraybackslash}X@{}}}%
  {\endtabularx\par}
\newcommand{\rowsep}{\\[2pt]}

\begin{document}
\small

% ---------------------------------------------------------------- header
\begin{center}
  {\sffamily\bfseries\fontsize{23}{27}\selectfont Minkyu Jeon}\\[4pt]
  {\sffamily\color{accent}Agentic AI for De Novo Protein Design
   \enspace$\cdot$\enspace Heterogeneous Cryo-EM Reconstruction
   \enspace$\cdot$\enspace Generative Models}\\[3pt]
  Ph.D. Candidate, Computer Science, Princeton University\\[2pt]
  \href{mailto:mj7341@princeton.edu}{mj7341@princeton.edu}\sep
  \href{https://minkyujeon.github.io}{minkyujeon.github.io}\sep
  \href{https://scholar.google.com/citations?hl=en&user=kyQR6ecAAAAJ}{Google Scholar}\sep
  \href{https://github.com/minkyujeon}{GitHub}\sep
  \href{https://www.linkedin.com/in/minkyu-jeon-9a5b92152/}{LinkedIn}\sep
  \href{https://x.com/MinkyuJeon19791}{X}
\end{center}

% ---------------------------------------------------------------- summary
\cvsection{Summary}
Ph.D. candidate developing generative and representation-learning methods for scientific
inverse problems, focused on (i) \textbf{agentic inference-time scaling for de novo protein
design}, where LLM agents run binder-design campaigns under a fixed GPU budget by scheduling
generative and verification tools (T-REX), and (ii) \textbf{heterogeneous cryo-EM
reconstruction} of conformational and compositional variability (CryoBench, CryoHype).
Research experience at Prescient Design (Genentech) and the Broad Institute.

% ---------------------------------------------------------------- education
\cvsection{Education}
\entry{Princeton University}{2023--present}
  {Ph.D. Candidate in Computer Science \enspace(advisor: Prof.\ Ellen D.\ Zhong)}{Princeton, NJ}
\entry{Korea University}{2020--2023}
  {M.S. in Computer Science and Engineering}{Seoul, Korea}
\entry{Konkuk University}{2014--2020}
  {B.S. in Applied Statistics and Computer Engineering (double major)}{Seoul, Korea}

% ---------------------------------------------------------------- research
\cvsection{Selected Research}
\entry{T-REX: Agentic Campaign Control for De Novo Binder Design}{2026--present}
  {First author \,$\cdot$\,
   \href{https://www.biorxiv.org/content/10.64898/2026.09.22.753604v2}{bioRxiv 2026}
   \,$\cdot$\, \href{https://github.com/ml-struct-bio/T-REX}{open-source code}}{}
\begin{itemize}
  \item LLM agents decide when to rescue, explore, or exploit, orchestrating six generation methods
        (Prote\'ina-Complexa $\times$4, BindCraft, BoltzGen), ProteinMPNN, and AlphaFold2
        evaluation; a deterministic controller validates and schedules every job.
  \item Highest throughput of structurally distinct in-silico hits on all 7 targets (1,461 total):
        \textbf{2.43$\times$} a tree-search controller and \textbf{2.48$\times$} the strongest
        generator-only baseline per target (geometric mean, matched worker-GPU budgets).
\end{itemize}
\vspace{2pt}

\entry{T-REX for Wet-Lab Binder Design Challenges}{2026}
  {Anthropic $\times$ Adaptyv 2026 \,$\cdot$\, EGFR, TNF-$\alpha$, GLP1R, TfR1, KRAS G12D pMHC}{}
% TODO: add wet-lab outcomes / rankings once Adaptyv results are released.
\begin{itemize}
  \item Added in-silico proxies for pH-selective binding (Proton-PottsMPNN redesign) and mouse
        cross-reactivity on EGFR, TNF-$\alpha$, and TfR1, plus structural-novelty diagnostics.
  \item Expanded the method pool (BindCraft2, PXDesign, RFdiffusion3, ArcRefine) and ran
        campaigns with Claude or Qwen as the planning LLM.
\end{itemize}
\vspace{2pt}

\entryone{Heterogeneous Cryo-EM Reconstruction \& Protein Conformation}{2023--present}
\begin{itemize}
  \item \textbf{CryoBench} (first author, NeurIPS 2024 Datasets \& Benchmarks \award{Spotlight}): benchmark
        datasets and metrics for conformational and compositional heterogeneity.
  \item \textbf{CryoHype} (co-first, CVPR 2026): transformer hypernetwork that reconstructs
        1,000 distinct structures from unlabeled images (fixed-pose).
  \item \textbf{Self-distillation for amortized reconstruction} (co-first, NeurIPS 2025 MLSB
        Workshop \award{Oral}).
  \item \textbf{ConforNets} (co-author): latent-space control of alternate conformations
        in OpenFold3.
\end{itemize}

% ---------------------------------------------------------------- experience
\cvsection{Experience}
\entry{Prescient Design, Genentech}{Nov 2025--Mar 2026}{Contractor (part-time)}{New York, NY}
\role{Research Intern}{Jun 2025--Aug 2025}
\begin{itemize}
  \item Led a project on understanding and improving 3D generative models (diffusion,
        flow matching) for protein design; studied limitations of voxel- and structure-space
        models and developed design and evaluation workflows.
\end{itemize}
\vspace{2pt}

\entry{Broad Institute of MIT and Harvard}{Jan 2023--Aug 2023}
  {Associate Computational Biologist \enspace(advisor: Dr.\ Hyunghoon Cho)}{Cambridge, MA}
\role{Visiting Graduate Student}{Sep 2021--Apr 2022}
\begin{itemize}
  \item Led privacy-preserving retinal image synthesis with GANs (\textbf{k-SALSA},
        first author, ECCV 2022) and later with diffusion models.
\end{itemize}
\vspace{2pt}

\entryone{Earlier Research Internships}{2018--2020, Seoul, Korea}
\begin{itemize}
  \item \textbf{Korea University} (Prof.\ Hyunwoo J.\ Kim), 2020: meta-learning for
        layer-wise optimization (IEEE Access 2021).
  \item \textbf{MakinaRocks}, 2019: anomaly detection with variational autoencoders.
  \item \textbf{KIST} (Dr.\ Yong Moo Kwon), 2018: situation-aware chatbot.
\end{itemize}

% ---------------------------------------------------------------- publications
\newpage
\cvsection{Publications}
{\color{muted}\textsuperscript{*}\,equal contribution}
\vspace{3pt}

{\sffamily\bfseries AI for Protein Design \& Structural Biology}\par\vspace{3pt}

\paper{bioRxiv 2026}{}{Agentic campaign control for high-throughput de novo binder design}
  {\me, Jinyeop Song, Jina Kim, Ellen D.\ Zhong
   \links{\href{https://www.biorxiv.org/content/10.64898/2026.09.22.753604v2}{paper} /
   \href{https://github.com/ml-struct-bio/T-REX}{code}}}

\paper{CVPR 2026}{}{CryoHype: Reconstructing a thousand cryo-EM structures with transformer-based hypernetworks}
  {Jeffrey Gu\eq, \me\eq, Ambri Ma, Serena Yeung-Levy, Ellen D.\ Zhong
   \links{\href{https://arxiv.org/abs/2512.06332}{paper}
   \enspace{\color{muted}also NeurIPS 2025 MLSB Workshop (poster)}}}

\paper{arXiv 2026}{}{ConforNets: Latents-based conformational control in OpenFold3}
  {Minji Lee, Colin Kalicki, \me, Aymen Qabel, Alisia Fadini, Mohammed AlQuraishi
   \links{\href{https://arxiv.org/abs/2604.18559}{paper}}}

\paper{NeurIPS 2025 MLSB Workshop}{Oral}{Separating signal from noise: A self-distillation approach for amortized heterogeneous cryo-EM reconstruction}
  {\me\eq, Jeffrey Gu\eq, Ambri Ma, Serena Yeung-Levy, Vincent Sitzmann, Ellen D.\ Zhong
   \links{\href{https://openreview.net/pdf?id=YIvDrdPCo2}{paper}}}

\paper{NeurIPS 2024 D\&B}{Spotlight}{CryoBench: Diverse and challenging datasets for the heterogeneity problem in cryo-EM}
  {\me, Rishwanth Raghu, Miro Astore, Geoffrey Woollard, Ryan Feathers, Alkin Kaz,
   Sonya M.\ Hanson, Pilar Cossio, Ellen D.\ Zhong
   \links{\href{https://arxiv.org/abs/2408.05526}{paper} /
   \href{https://cryobench.cs.princeton.edu/}{project}}}

\vspace{1pt}
{\sffamily\bfseries Machine Learning, Computer Vision \& Health}\par\vspace{3pt}

\paper{ECCV 2026 Workshop (MUCG)}{}{CHIMERA: Adaptive cache injection and semantic anchor prompting for zero-shot image morphing with morphing-oriented metrics}
  {Dahyeon Kye, Jeahun Sung, \me, Jihyong Oh
   \links{\href{https://www.arxiv.org/abs/2512.07155}{paper} /
   \href{https://cmlab-korea.github.io/CHIMERA/}{project}}}

\paper{arXiv 2025}{}{R3eVision: A survey on robust rendering, restoration, and enhancement for 3D low-level vision}
  {Weeyoung Kwon, Jeahun Sung, \me, Chanho Eom, Jihyong Oh
   \links{\href{https://arxiv.org/abs/2506.16262}{paper}}}

\paper{JMIR Med.\ Inform.\ 2024}{}{A multivariable prediction model for mild cognitive impairment and dementia: Algorithm development and validation}
  {Sarah Soyeon Oh, Bada Kang, Dahye Hong, Jennifer Ivy Kim, Hyewon Jeong, Jinyeop Song, \me
   \links{\href{https://doi.org/10.2196/59396}{paper}}}

\paper{Inf.\ Sci.\ 2023}{}{Randomly shuffled convolution for self-supervised representation learning}
  {Youngjin Oh\eq, \me\eq, Dohwan Ko, Hyunwoo J.\ Kim
   \links{\href{https://www.sciencedirect.com/science/article/pii/S0020025522013032}{paper} /
   \href{https://github.com/minkyujeon/CROFFLE/tree/main/}{code}}}

\paper{NeurIPS 2022}{}{SageMix: Saliency-guided mixup for point clouds}
  {Sanghyeok Lee\eq, \me\eq, Injae Kim, Yunyang Xiong, Hyunwoo J.\ Kim
   \links{\href{https://openreview.net/pdf?id=q-FRENiEP_d}{paper} /
   \href{https://github.com/mlvlab/SageMix/}{code}}}

\paper{ECCV 2022}{}{$k$-SALSA: $k$-anonymous synthetic averaging of retinal images via local style alignment}
  {\me, Hyeonjin Park, Hyunwoo J.\ Kim, Michael Morley, Hyunghoon Cho
   \links{\href{https://arxiv.org/abs/2303.10824}{paper} /
   \href{https://github.com/hcholab/k-salsa/}{code}}}

\paper{IEEE Access 2021}{}{Learning to balance local losses via meta-learning}
  {Seungdong Yoa, \me, Youngjin Oh, Hyunwoo J.\ Kim
   \links{\href{https://ieeexplore.ieee.org/document/9541196}{paper}}}

% ---------------------------------------------------------------- talks
\cvsection{Talks}
\begin{rows}
  Dec 2025 & NeurIPS 2025 MLSB Workshop (oral, with Jeffrey Gu): amortized heterogeneous cryo-EM reconstruction\rowsep
  Aug 2024 & Flatiron Institute, Simons Foundation (invited): CryoBench\rowsep
  Feb 2024 & CMU MetaMobility Lab (invited lecture): $k$-SALSA
\end{rows}

% ---------------------------------------------------------------- honors
\cvsection{Honors}
\begin{rows}
  2024--2026 & \textbf{Asan Foundation Biomedical Science Scholarship} (three-year doctoral scholarship)\rowsep
  2021 & \textbf{I-Corps Program}, \$35,000 government support for startup customer discovery\rowsep
  2019 & Dean's List and 40\% tuition scholarship, Konkuk University\rowsep
  2016 & Certificate of Commendation, Republic of Korea Marine Corps (military service 2015--2016)
\end{rows}

% ---------------------------------------------------------------- service & teaching
\cvsection{Service \& Teaching}
\begin{rows}
  Reviewing & AAAI 2027 (program committee); NeurIPS 2025--26; ICLR 2025--27; CVPR 2026;
             ECCV 2026; NeurIPS workshops 2023--24; \textit{Information Sciences} 2024\rowsep
  Princeton & TA, COS302 Mathematics for Numerical Computing and ML (Fall 2024);
              Instructor, AI4ALL (Summer 2024)\rowsep
  Korea Univ. & TA, Artificial Intelligence (2021) and AI Security (2020);
                Tutorial Teaching Fellow, Deep Learning (2020)
\end{rows}

% ---------------------------------------------------------------- skills
\cvsection{Skills}
\begin{rows}
  Machine learning & PyTorch, TensorFlow; diffusion and flow matching; transformers;
                     implicit neural representations; LLM agents (vLLM)\rowsep
  Computing & Python, R, C; Linux; SLURM multi-GPU clusters; Git
\end{rows}

\end{document}
